\chapter{Dokumentation zu Hoover}
\label{DokumentationHooverCrawler}

\section{Starten der Fetcher}
Wie schon ausf"uhrlich er"ortert, laufen die Fetcher, die die eigentliche Arbeit des {\em Hoover}-Systems erledigen, auf vielen Rechnern. Auf jedem, dem System zuarbeitenden Rechner, mu"s also ein Fetcher laufen, der sich dann beim {\em Hoover} anmeldet.
Der eigentliche Programmaufruf sieht wie folgt aus:
\begin{center}
{\tt \small Fetcher [-threads \flq number\frq ] [-host \flq hostname\frq ]}
\end{center}
Die Parameter sind optional. In der Voreinstellung wird mit f"unfzig Threads gefetched und der lokale Rechner als {\em Hoover} angenommen.


\section{Starten von Hoover}

Hoover ist, wie schon in Kapitel \ref{Crawler} besprochen, mit einem Konfigurationsfile zu starten. Der eigentliche Programmaufruf sieht dann wie folgt aus:
\begin{center}
{\tt \small Hoover [ConfigurationFileName]}
\end{center}
Der Parameter ist optional. In der Voreinstellung wird mit "`HooverConfiguration"' als Konfigurationsfile gearbeitet. Hoover gibt die erhaltenen HTML-Dokumente auf die Standardausgabe aus, so da"s der Programmaufruf in einer Unix-Shell meist wie folgt aussieht:
\begin{center}
{\tt \small Hoover \frq fetchedWebDokuments 2\frq hooverProgressInformation }
\end{center}
Im Klartext: Alle geholten HTML-Dokumente werden in die Datei "`fetchedWebDokuments"' geschrieben, die Informationen, was {\em Hoover} gerade macht, in die Datei "`hooverProgressInformation"'.


\subsection{Das Hoover Konfigurationsfile}

Das Konfigurationsfile von Hoover entspricht dem "'PropertyList"'-Format von OpenStep. Dieses Format erlaubt es Daten aus Textfiles einzulesen und in Objekte abzubilden.

Ein {\em NSDictionary}-Objekt erzeugt man durch Eingabe von "`Schl"usselName = WertName"', ein {\em NSString}-Objekt wird von G"ansef"uschen umramt, ein {\em NSArray} erzeugt man, indem man die Werte mit Komma trennt und in runde Klammern setzt. Kommentare werden zwischen "`/*"' und "`*/"' geschrieben.

Soviel zur Theorie, in der Praxis sieht ein HooverKonfigurationsfile z.B. so aus:

{\small
\begin{verbatim}
{   /*
        Dies ist eine TestKonfiguration, um zu zeigen wie einfach es
        ist ein poropertyList Format von Hand einzugeben
    */
    urls =  (   "http://www.leo.org/demap/cities/Muenchen.html",
                "http://watzmann.cis.uni-muenchen.de/Deutschland.html"
            );
    general = {
            useragentname = "Hoover/0.1";
            useragentmail = "Patrick Stein <jolly@cis.uni-muenchen.de>";
            databasename =  "webdatabase" ;
            includehosts = ( ".de",".at","ch" );
            excludehosts = ();
            includepaths = ();
            excludepaths = ( "?", "cgi-bin", "cgi_bin" );
            httpproxy = { host="www.cis.uni-muenchen.de"; port = 8000 };
        };
}
\end{verbatim}
}

Die Datei enth"alt zum gegenw"artigen Zeitpunkt ein W"orterbuch mit den folgenden zwei Eintr"agen:
\begin{description}
	\item[urls:]	Hier kann man die Startseiten, mit denen {\em Hoover} beginnen soll, als {\em URL}-Array eingeben.
	\item[general:]	Hier werden Schl"ussel-Wert Paare aufgelistet, die die verschiedenen Optionen im {\em Hoover} beeinflussen:
	\begin{description}
		\item[useragentname:]	Mit diesem Namen meldet sich {\em Hoover} bei jeder Web-Site. Diese k"onnen dann den Zugriff "uber das {\em robots exclusion protocol} mit diesem Namen verweigern oder eingrenzen.
		\item[useragentmail:]	Mite dieser E-Mail Adresse meldet sich {\em Hoover} bei jeder Web-Site. So k"onnen die Administratoren bei Problemen mit {\em Hoover} dies gleich der verantwortlichen Person mitteilen.
		\item[databasename:]	Unter diesem Dateinamen wird die persistente Datenspeicherung der {\em URLs} durchgef"uhrt.
		\item[includehosts:]	Es werden Hosts nur von {\em Hoover} kontaktiert, wenn mindestens eine Zeichenkette in diesem Array auf den Hostnamen pa"st ( matched ).
		\item[excludehosts:]	Hosts, deren Hostname auf mindestens eine Zeichenkette in diesem Array passt, werden niemals kontaktiert.
		\item[includepaths:]	Siehe includehosts, nur das hier auf den Pfadnamen gematched wird.
		\item[excludepaths:]	Siehe excludehosts, nur das hier auf den Pfadnamen gematched wird.
		\item[httpproxy:]		Existiert dieser Eintrag, dann wird das Internet von den Fetchern nur "uber den benannten Host (HTTP-Proxy) kontaktiert. Bei Nichtexistenz wird das Internet direkt von den Fetchern benutzt.
	\end{description}
\end{description}

Es hat sich w"ahrend der Implementierung der Vorteil dieser {\em variablen} Konfiguration gezeigt. Der Eintrag "`httpproxy"' war nicht von Anfang an in der Konfiguration enthalten.

Der Eintrag wurde notwendig, als das {\em Hoover}-System "uber einen Firewall-Rechner auf das Internet zugreifen sollte. Ich erg"anzte daraufhin die Spezifikation und Implementation und konnte nun das System "uber einen Firewall auf das Internet zugreifen lassen.

Bahnbrechend ist dies sicher nicht, jedoch ist es ohne Probleme m"oglich die neuen, erweiterten Konfigurationsfiles auf den {\em alten} {\em Hoovern} laufen zu lassen, da diese Systeme gar nicht nach einem "`httpproxy"'-Eintrag im "'general"'-W"orterbuch suchen.


